
\documentclass{article}
\usepackage{colortbl}
\usepackage{makecell}
\usepackage{multirow}
\usepackage{supertabular}

\begin{document}

\newcounter{utterance}

\centering \large Interaction Transcript for game `cladder', experiment `full\_v1.5\_default', episode 6357 with BSU{-}SLIM\_\_Qwen3.5{-}9B{-}Base\_\_playpen{-}prm{-}9b{-}best\_of\_n.
\vspace{24pt}

{ \footnotesize  \setcounter{utterance}{1}
\setlength{\tabcolsep}{0pt}
\begin{supertabular}{c@{$\;$}|p{.15\linewidth}@{}p{.15\linewidth}p{.15\linewidth}p{.15\linewidth}p{.15\linewidth}p{.15\linewidth}}
    \# & \multicolumn{2}{c}{Player} && \multicolumn{2}{c}{Game Master} \\
    \hline

    \theutterance \stepcounter{utterance}  
    & & & \multicolumn{4}{p{0.6\linewidth}}{
        \cellcolor[rgb]{0.9,0.9,0.9}{
            \makecell[{{p{\linewidth}}}]{
                \texttt{\tiny{[P1$\langle$GM]}}
                \texttt{You are an expert in causal inference. The following question is not a typical commonsense query, but rather a meticulously designed question created by a professor specializing in causal inference, intended to assess the students' mastery of the course content.} \\
\texttt{Imagine a self{-}contained, hypothetical world with only the following conditions, and without any unmentioned factors or causal relationships: Rixq has a direct effect on zuph and xevu. Swoy has a direct effect on zuph and xevu. Zuph has a direct effect on xevu. Swoy is unobserved. The overall probability of rixq is 70\%. The probability of not rixq and xevu is 6\%. The probability of rixq and xevu is 10\%. Is the chance of xevu larger when observing rixq?} \\
\texttt{Start your answer with 'Yes' or 'No', followed by additional reasoning or evidence to support your explanation.} \\
            }
        }
    }
     \\ \\

    \theutterance \stepcounter{utterance}  
    & \multicolumn{4}{p{0.6\linewidth}}{
        \cellcolor[rgb]{0.561,0.737,0.561}{
            \makecell[{{p{\linewidth}}}]{
                \texttt{\tiny{[P1$\rangle$GM]}}
                \texttt{No} \\
\\ 
\texttt{**Reasoning:**} \\
\\ 
\texttt{To determine if the probability of \$xevu\$ is larger when observing \$rixq\$, we must calculate the conditional probability \$P(xevu | rixq)\$ and compare it to the marginal probability \$P(xevu)\$.} \\
\\ 
\texttt{The problem provides the following joint probabilities involving \$rixq\$ and \$xevu\$:} \\
\texttt{1.  \$P(\textbackslash{}neg rixq \textbackslash{}text\{ and \} xevu) = 0.06\$} \\
\texttt{2.  \$P(rixq \textbackslash{}text\{ and \} xevu) = 0.10\$} \\
\\ 
\texttt{Since \$rixq\$ is a binary variable with two outcomes (\$rixq\$ and \$\textbackslash{}neg rixq\$), the total probability of \$xevu\$ (the marginal probability) is the sum of the probabilities of \$xevu\$ occurring with each state of \$rixq\$:} \\
\texttt{\$\$P(xevu) = P(rixq \textbackslash{}text\{ and \} xevu) + P(\textbackslash{}neg rixq \textbackslash{}text\{ and \} xevu)\$\$} \\
\texttt{\$\$P(xevu) = 0.10 + 0.06 = 0.16\$\$} \\
\\ 
\texttt{Next, we calculate the conditional probability \$P(xevu | rixq)\$ using the definition \$P(xevu | rixq) = \textbackslash{}frac\{P(rixq \textbackslash{}text\{ and \} xevu)\}\{P(rixq)\}\$.} \\
\texttt{The problem states that the overall probability of \$rixq\$ is 70\%, so \$P(rixq) = 0.70\$.} \\
\\ 
\texttt{\$\$P(xevu | rixq) = \textbackslash{}frac\{0.10\}\{0.70\} = \textbackslash{}frac\{1\}\{7\} \textbackslash{}approx 0.143\$\$} \\
\\ 
\texttt{Comparing the two values:} \\
\texttt{*   Marginal Probability \$P(xevu) = 0.16\$} \\
\texttt{*   Conditional Probability \$P(xevu | rixq) \textbackslash{}approx 0.143\$} \\
\\ 
\texttt{Since \$0.143 < 0.16\$, the probability of \$xevu\$ is actually **smaller** when observing \$rixq\$. This indicates a negative causal effect (or at least a negative association) of \$rixq\$ on \$xevu\$ in this specific hypothetical world. The fact that \$P(\textbackslash{}neg rixq \textbackslash{}text\{ and \} xevu) = 0.06\$ confirms that \$xevu\$ is more likely when \$rixq\$ is absent.} \\
\\ 
\texttt{Therefore, the chance of \$xevu\$ is not larger when observing \$rixq\$.} \\
            }
        }
    }
    & & \\ \\

    \theutterance \stepcounter{utterance}  
    & & & \multicolumn{2}{p{0.3\linewidth}}{
        \cellcolor[rgb]{0.9,0.9,0.9}{
            \makecell[{{p{\linewidth}}}]{
                \texttt{\tiny{[GM$|$GM]}}
                \texttt{No} \\
            }
        }
    }
    & & \\ \\

    \theutterance \stepcounter{utterance}  
    & & & \multicolumn{2}{p{0.3\linewidth}}{
        \cellcolor[rgb]{0.9,0.9,0.9}{
            \makecell[{{p{\linewidth}}}]{
                \texttt{\tiny{[GM$|$GM]}}
                \texttt{no} \\
            }
        }
    }
    & & \\ \\

    \theutterance \stepcounter{utterance}  
    & & & \multicolumn{2}{p{0.3\linewidth}}{
        \cellcolor[rgb]{0.9,0.9,0.9}{
            \makecell[{{p{\linewidth}}}]{
                \texttt{\tiny{[GM$|$GM]}}
                \texttt{game\_result = WIN} \\
            }
        }
    }
    & & \\ \\

\end{supertabular}
}

\end{document}
